\subsection{Probabilistic no-signaling projections}

\begin{proposition}[Probabilistic finite empirical model]\label{thm:probabilistic-model}\lean{Obsmetrologia.QuantumContextuality.PModel}\leanok\uses{thm:measurement-scenario}
A probabilistic empirical model assigns a finite weighted distribution to each
context.
\end{proposition}

\begin{theorem}[A global distribution is no-signaling]\label{thm:global-dist-no-signaling}\lean{Obsmetrologia.QuantumContextuality.globalDist_implies_noSignaling}\leanok\uses{thm:probabilistic-model}
Any model induced by one global weighted distribution satisfies no-signaling:
shared-measurement marginals agree across overlapping contexts.
\end{theorem}

\begin{theorem}[The PR distribution is no-signaling]\label{thm:pr-box-p-no-signaling}\lean{Obsmetrologia.QuantumContextuality.prBoxP_noSignaling}\leanok\uses{thm:probabilistic-model}
The finite PR distribution is no-signaling even though its support is
contextual.
\end{theorem}

\begin{theorem}[The no-signaling PR support is strongly contextual]\label{thm:pr-support-strong-contextual}\lean{Obsmetrologia.QuantumContextuality.prBoxP_support_strongly_contextual}\leanok\uses{thm:pr-box-p-no-signaling, thm:pr-box-strong-contextual}
Forgetting weights from the no-signaling PR distribution yields a support
model with no global section.
\end{theorem}
